%\subsection{An End-to-End Decision Tree Benchmark}\label{eva:pdte}

% After query process, we use the function described in \cref{Analysis: Bridge} to bridge the gap between 
% aggregation result and data analysis. 
% {\bf{End-to-End Decision Tree Benchmark}}: 
%\paragraph{End-to-End Decision Tree Benchmark}
For the evaluation of the private decision tree, 
we provide a brief summary of the evaluation results obtained over
multiple decision tree models under a computation outsourcing setting, where both
the inputs and the model weights of the decision trees are ciphertexts.
% through our \pdte on 
% the experiment environment described \cref{eva:Arch}. 
To enable fair comparisons, we use the same datasets 
as in~\cite{sortinghat}, and reworked the implementation of~\cite{sortinghat} such that the
solution in~\cite{sortinghat} can also operate over encrypted inputs and model weights.
%\cref{fig: homtreebench} compares the performance of \NAME and \cite{sortinghat} on  
%end-to-end decision tree benchmark. 
While detailed comparison results can be found in the appendix, we achieve 12x--326x faster 
evaluation speed than~\cite{sortinghat} over a wide range of datasets and input sizes.
Thus, by adding the powerful online data analysis
capability to the expressive query evaluation framework,
\NAME is able to significantly improve the efficiency, usability
and practical applicability of outsourcing databases securely.
% able to significantly narrow
% the performance and usability gaps between theoretical FHE-based encrypted
% database constructions and practical needs for SDO in the real world.
% improve the 
% practicality of secure database outsourcing based on fully homomorphic encryption.
